\chapter{Evaluación}
\label{cap:evaluacion}

Este capítulo presenta los resultados empíricos del sistema SGROAS organizados por
bloque de evaluación (Bloque C de la guía \cite{uteq2026guia}). Cada bloque declara:
la pregunta de investigación que responde, la métrica empleada, los datos crudos que
la soportan, la estadística descriptiva con intervalo de confianza al 95\,\% y la
interpretación. Todos los números son generados con scripts versionados y pueden
reproducirse siguiendo \texttt{docs/mediciones/DATA-PROVENANCE.md}. El diseño
(GQM, RQs y umbrales declarados \emph{a priori}) está en
\texttt{docs/estudio/DISENO-ESTUDIO.md}.

\section{Rendimiento (Bloque C.1) — RQ1}
\label{sec:perf}

\textbf{Pregunta:} \emph{¿El acceso híbrido a datos CRUD/ORM + stored procedures
mantiene un tiempo de respuesta aceptable bajo carga?}

\textbf{Métricas:} latencia de \texttt{http\_req\_duration} (media, mediana,
percentil 95) y tasa de error \texttt{http\_req\_failed} sobre
\texttt{GET /api/conductores}, con 50 VUs durante 30\,s (\texttt{k6/opts.js}).
El umbral declarado es p95~<~200\,ms y error~<~1\,\%.

\textbf{Datos crudos:} tres corridas independientes en
\texttt{docs/mediciones/perf/k01-run1.json}, \texttt{k02-run2.json} y
\texttt{k03-run3.json} (commit \texttt{62bf8fa}).

\textbf{Descriptiva e IC 95\,\%:} la media de las medias por corrida es
\textbf{23.01\,ms} (DT 12.44\,ms, EE 7.18\,ms), con IC\,95\,\%
\textbf{[$-7.88$; $53.91$]\,ms} (t de Student, gl = 2, n = 3 corridas). El p95
máximo es 173.13\,ms (corrida k01), por debajo del umbral de 200\,ms, y la tasa de
error es 0\,\% con 8930 verificaciones correctas. La primera corrida (37.3\,ms de
media; p95 173.13\,ms) refleja el arranque en frío de conexiones y del cache
Redis; las corridas k02 y k03 estabilizan la media por debajo de 17\,ms
(\texttt{docs/mediciones/perf/ANALISIS-k6.md}).

\textbf{Interpretación:} el sistema cumple los umbrales de rendimiento en la
condición de cache caliente. Dado que las tres corridas corresponden a la misma
condición, el contraste formal cache caliente-frío queda definido a priori en
\texttt{scripts/perf/nonparametric.py} (U de Mann-Whitney, Wilcoxon y d de Cliff
\cite{mann_whitney1947,wilcoxon1945,cliff1993}) y se ejecutará con la condición de
cache fría cuando la URL pública esté disponible (K3). La medición sigue los
lineamientos de pruebas de carga de sistemas de gran escala de \cite{jiang2015}.

\section{Seguridad (Bloque C.2) — RQ2}
\label{sec:sec}

\textbf{Pregunta:} \emph{¿El sistema mitiga los riesgos más críticos del OWASP Top
10 \cite{owasp2021} con controles verificables?}

\textbf{Evidencia (datos crudos):} se documentan seis controles en
\texttt{docs/mediciones/sec/}: control de acceso (A01), criptografía de contraseñas
y sesión (A02), prevención de inyección con repositorios Spring Data y
\texttt{@Procedure} sin SQL concatenado (A03), cabeceras de seguridad y error
\emph{Problem Details} RFC\,7807 \cite{rfc7807} (A05), limitación de tasa de login
(A07) y registros de auditoría (A09).

\textbf{Medidas mecánicas:} la regla de oro del proyecto prohíbe SQL dinámico
concatenado; se audita con \texttt{scripts/audit-sql-dynamic.sh} (commit
\texttt{1a07dc7}) y con SpotBugs/Find-Sec-Bugs en CI. El esquema de autenticación
usa JWT según RFC\,7519 \cite{rfc7519} sobre OAuth\,2.0 \cite{rfc6749}, con cookie
HttpOnly y BCrypt (factor 10+).

\textbf{ZAP (pendiente de URL):} el escaneo baseline de OWASP ZAP
(\texttt{scripts/zap/run-zap.sh}) se ejecutará contra la URL pública y su resultado
se reportará en \texttt{docs/mediciones/sec/zap/RESUMEN.md} (dependencia K3).

\textbf{Interpretación:} los controles estáticos y dinámicos ejecutados no detectan
SQL dinámico ni credenciales en texto plano; la evidencia A01 (accessible via
token) confirma el cierre del acceso no autenticado.

\section{Usabilidad (Bloque C.3) — RQ3}
\label{sec:sus}

\textbf{Pregunta:} \emph{¿Qué nivel de usabilidad perciben los usuarios del
sistema?}

\textbf{Instrumento:} System Usability Scale de Brooke \cite{brooke1996}, 10 ítems
Likert 1--5. Puntuación = (suma de contribuciones) $\times$ 2.5.

\textbf{Datos crudos:} 10 participantes (P01--P10) anonimizados en
\texttt{docs/mediciones/sus/sus-raw.csv} y por participante en
\texttt{P01.json..P10.json}; consentimientos custodiados fuera del repositorio.

\textbf{Descriptiva e IC 95\,\%:} media \textbf{63.0/100}, DT 13.88, EE 4.39, IC\,95\,\%
\textbf{[53.07; 72.93]} (t = 2.262, gl = 9), mínimo 47.5 y máximo 90.0
(\texttt{docs/mediciones/sus/ANALISIS-SUS.md}). En la escala adjetiva de Bangor et
al. \cite{bangor2009,bangor2008} la media corresponde a \emph{Bueno}, dentro de la
zona de aceptabilidad \emph{marginal} (50--70) \cite{lewis2009}.

\textbf{Interpretación:} la media no alcanza el umbral de 70 que Bangor asocia a
sistemas aceptables; se documenta como área de mejora prioritaria (ítems pares de
complejidad percibida). El tamaño muestral (n = 10) cumple la recomendación mínima
para SUS \cite{kortum2013}.

\section{Cobertura de pruebas (Bloque C.4)}
\label{sec:cobertura}

\textbf{Métrica:} contadores de JaCoCo (instrucciones, ramas y líneas) sobre la
ejecución de \texttt{./mvnw verify} con 150 pruebas JUnit 5.

\textbf{Datos crudos:} \texttt{docs/mediciones/jacoco/jacoco.csv} y reporte HTML
(commit \texttt{5f3b472}).

\textbf{Resultado:} instrucciones 98.8\,\%, ramas 85.4\,\% y líneas 99.7\,\%, con
0 fallos. Se supera el umbral del proyecto (líneas $\ge$ 0.70 para la entrega
final). La discusión sobre la relación entre cobertura y detección de fallos se
fundamenta en \cite{inozemtseva2014,gopinath2014,gligoric2013,ivankovic2019}.

\section{Calidad web (Bloque C.5) — RQ4}
\label{sec:lighthouse}

\textbf{Pregunta:} \emph{¿La interfaz web cumple estándares de rendimiento,
accesibilidad, mejores prácticas y SEO?}

\textbf{Métrica:} categorías de Lighthouse \cite{jacoco} (móvil, throttling Slow 4G)
sobre el login del frontend Angular.

\textbf{Datos crudos:} \texttt{docs/mediciones/lighthouse/lhci-20260730-2115.json} y
\texttt{...2117.json} (commit \texttt{1a07dc7}).

\textbf{Resultado:} Performance 100, Accessibility 95, Best Practices 100, SEO 90;
los cuatro umbrales (80/90/90/90, \texttt{lighthouserc.js}) se cumplen. Las
corridas adicionales m1--m3 y d1--d3 se ejecutarán con la URL pública
(\texttt{scripts/lighthouse/run-lighthouse.sh}).

\textbf{Interpretación:} el SPA cumple los estándares; SEO (90) es el punto más
cercano al umbral y se documenta como mejora menor (metadatos).

\section{Síntesis}
\label{sec:sintesis}

La Tabla~\ref{tab:sintesis} resume el cumplimiento por bloque.

\begin{table}[H]
\centering
\caption{Síntesis de resultados por bloque de evaluación.}
\label{tab:sintesis}
\begin{tabular}{@{}llll@{}}
\toprule
Bloque & Indicador & Valor & Cumple \\
\midrule
C.1 Rendimiento & p95 $<$ 200\,ms (cache caliente) & 173.13\,ms & Sí \\
C.1 Rendimiento & error rate $<$ 1\,\% & 0\,\% & Sí \\
C.2 Seguridad & SQL dinámico / A01--A09 & 0 violaciones & Sí \\
C.3 Usabilidad & SUS $\ge$ 70 & 63.0 & No (marginal) \\
C.4 Cobertura & líneas $\ge$ 70\,\% & 99.7\,\% & Sí \\
C.5 Calidad web & Perf/Acc/BP/SEO $\ge$ 80/90/90/90 & 100/95/100/90 & Sí \\
\bottomrule
\end{tabular}
\end{table}